\chapter{Information Geometry and Scale Thermodynamics}

Chapter 2 showed that the holographic boundary is not a periodic crystal but an aperiodic symbolic system generated by forbidden fivefold order.  We now equip this boundary with thermodynamics.  The central claim of this chapter is that the flow from bulk to boundary is an information-geometric scaling process whose terminal states are nilpotent zero modes.  In algebraic form, the thermodynamic endpoint is not a large-energy state but a square-zero collapse:
\begin{equation}
    Z^2=0.
\end{equation}
These nilpotent modes are interpreted as massless parafermionic boundary sinks: they absorb scale flow while carrying only topological entropy.

\section{Modular Flow as Scale Thermodynamics}

The bulk/boundary transition is governed by a modular scaling generator.  In the biquaternionic bulk, boosts and scale transformations are represented by exponentials of traceless Pauli-sector operators.  If $K$ is such a modular generator, then the flow takes the schematic form
\begin{equation}
    \rho_s = e^{-sK}\rho_0 e^{sK}.
\end{equation}
The parameter $s$ is not merely time; it is holographic depth.  Increasing $s$ strips away invertible bulk degrees of freedom and reveals the nilpotent boundary radical.

In the formal stack, this layer is represented by the modules
\begin{center}
\texttt{ModularItakuraBiquaternion.lean},\quad
\texttt{ModularHolographicMetric.lean},\quad
\texttt{NilpotentItakuraSaito.lean}.
\end{center}
The Lean proofs verify the algebraic kernels of the modular identities, while analytic convergence and physical interpretation are retained as theorem-honest sockets.

\section{The Itakura--Saito Divergence}

The scale functional used throughout the architecture is the noncommutative analogue of the Itakura--Saito divergence.  In scalar form, for positive variables $x,y>0$, it is
\begin{equation}
    D_{IS}(x\|y)=\frac{x}{y}-\log\left(\frac{x}{y}\right)-1.
\end{equation}
For matrix or operator variables, the logarithm and exponential are interpreted through functional calculus where available, and through algebraic series witnesses in the formal computational layer.

A useful normalized operator expression is
\begin{equation}
    D_{IS}(K)=e^K-K-I.
\end{equation}
When $K$ is nilpotent and square-zero, the exponential truncates:
\begin{equation}
    K^2=0
    \quad\Longrightarrow\quad
    e^K=I+K.
\end{equation}
Therefore
\begin{equation}
    D_{IS}(K)=0.
\end{equation}
This is the algebraic reason nilpotent zero modes become thermodynamic sinks: they carry no further scale divergence.

\begin{theorem}[Nilpotent Itakura--Saito Collapse]
If $K^2=0$, then the normalized Itakura--Saito defect vanishes:
\begin{equation}
    e^K-K-I=0.
\end{equation}
\end{theorem}

\begin{proof}[Formal certificate]
The square-zero truncation of the exponential is verified in the SymPy witness \texttt{nilpotent\_itakura\_saito.py}.  The Lean layer records the algebraic collapse by the theorem family in \texttt{NilpotentItakuraSaito.lean}, and the master certificate includes the square-zero statement as
\begin{equation}
    \texttt{Znil\_square\_zero}:\quad Z_{nil}Z_{nil}=0.
\end{equation}
\end{proof}

\section{Fisher/Bures Metric Closure}

Before the boundary collapse, the bulk still carries a smooth information metric.  The Fisher/Bures metric is the infinitesimal quadratic form measuring distinguishability of nearby states.  In the Pauli-biquaternion sector, a tangent vector has the form
\begin{equation}
    V=v_1\sigma_1+v_2\sigma_2+v_3\sigma_3.
\end{equation}
Using the Pauli relations,
\begin{equation}
    V^2=(v_1^2+v_2^2+v_3^2)I.
\end{equation}
Thus the metric closes isotropically in the bulk.  This is the same quadratic closure that appears in the universal polynomial symmetry operator formalism:
\begin{equation}
    (a\cdot\sigma)^2=(a\cdot a)I.
\end{equation}

This closure is verified in
\begin{center}
\texttt{PolynomialSymmetryOperators.lean}\quad via \quad\texttt{pauliVec\_sq},
\end{center}
and in SymPy by \texttt{polynomial\_symmetry\_operators.py}.

\section{Zero Modes and Thermal Sinks}

A nilpotent boundary operator
\begin{equation}
    Z=\begin{pmatrix}0&1\\0&0\end{pmatrix}
\end{equation}
satisfies
\begin{equation}
    Z^2=0.
\end{equation}
It has only the eigenvalue $0$, but it is not the zero operator.  This distinction is essential: the boundary zero mode is spectrally collapsed but algebraically active.  It can mediate extensions, defects, and parafermionic boundary transitions while contributing no ordinary bulk mass scale.

In thermodynamic language, the nilpotent sector has no quadratic energy storage.  It therefore behaves like a zero-specific-heat sink:
\begin{equation}
    C_v=0.
\end{equation}
The physical statement is encoded as a socket, while the algebraic cause is proved concretely as square-zero nilpotency.

\section{Majorana and Parafermionic Boundary Modes}

The superconducting wallpaper-fermion layer supplies a physical interpretation of these zero modes.  In the Bogoliubov--de Gennes setting, particle-hole symmetry pairs positive and negative energy modes.  At the nodal boundary, the zero-energy sector is naturally Majorana-like.  When nonsymmorphic glide and quasicrystalline boundary order are included, these Majorana modes organize into parafermionic and topological defect sectors.

The related formal modules include
\begin{center}
\texttt{WallpaperFermionSuperconductingGap.lean},\quad
\texttt{SuperchargeSquare.lean},\quad
\texttt{ProjectiveGlideSuperchargeUnification.lean}.
\end{center}
The key algebraic pattern is again polynomial:
\begin{equation}
    Q^2=H,
    \qquad
    q^2=0.
\end{equation}
The bulk supercharge squares to an even Hamiltonian, whereas the boundary degeneration squares to zero.

\section{Thermodynamic Collapse as Nilradical Projection}

Combining the above pieces, the holographic scale flow can be read as a projection from invertible bulk operators to the nilradical of the boundary algebra:
\begin{equation}
    \text{invertible bulk sector}
    \quad\longrightarrow\quad
    \text{zero-divisor boundary sector}
    \quad\longrightarrow\quad
    \text{nilpotent zero modes}.
\end{equation}

This is not a metaphor.  In the polynomial operator language, it is the degeneration of high-degree invertible relations into annihilating relations:
\begin{equation}
    A^{-1}\ \text{exists in the bulk},
    \qquad
    Z^2=0\ \text{at the boundary}.
\end{equation}
The boundary is therefore the algebraic place where scale thermodynamics terminates.

\section{Chapter Summary}

Chapter 3 identifies the thermodynamic meaning of the holographic boundary.  The Itakura--Saito scale divergence measures the loss of bulk distinguishability under modular flow.  The Fisher/Bures metric closes in the Pauli-biquaternion bulk, but the flow terminates at nilpotent boundary operators.  These square-zero modes are interpreted as Majorana/parafermionic zero modes and as thermodynamic sinks with vanishing ordinary heat capacity.

The chapter's central principle is
\begin{equation}
    \boxed{
    Z^2=0
    \quad\Longleftrightarrow\quad
    \text{zero-mode thermodynamic collapse}
    }
\end{equation}

Chapter 4 turns from thermodynamic collapse to spectral topology, showing that the same boundary fixed-set logic appears simultaneously in the Brillouin Klein bottle and the Riemann critical line.
